\begin{helper}
Let \(V,V'\), and \(\Lambda\) be finite types with decidable equality, and
let the source and target state spaces
\(\mathcal P(V)\times\mathbb R^V\) and
\(\mathcal P(V')\times\mathbb R^{V'}\) carry the source and target laws
\(\nu,\nu'\).  For each base index \(\omega_0\), let
\[
C^0_{\omega_0}\subseteq \mathcal P(V)\times\mathbb R^V,\qquad
C^{\prime0}_{\omega_0}\subseteq \mathcal P(V')\times\mathbb R^{V'}
\]
be paired supported base cells.  Let
\(Q_\lambda\) and \(Q'_\lambda\), \(\lambda\in\Lambda\), be the finite list
of source and target Boolean labels used by the Hybrid secondary refinement,
and let \(T_{\omega_0}\) be the paired transport relation on the base cells.
Assume the base cells and labels are measurable, and that \(T_{\omega_0}\)
transports membership in the base cells and in every label \(Q_\lambda\) to
the corresponding \(Q'_\lambda\).

For a truth vector \(\tau\subseteq\Lambda\), define the refined atoms
\[
C^0_{\omega_0,\tau}
=C^0_{\omega_0}\cap
  \{\eta:\eta\in Q_\lambda\Longleftrightarrow\lambda\in\tau
      \text{ for all }\lambda\in\Lambda\},
\]
and define \(C^{\prime0}_{\omega_0,\tau}\) analogously using the target
labels.  Assume that, for each \((\omega_0,\tau)\), the refined source and
target atoms already carry the common real mass certificate supplied by the
actual positive-mass product-subcell analysis on the supported common atom
family, together with the zero-mass branch: there is
\(w_{\omega_0,\tau}\ge0\) whose extended nonnegative value is both the source
atom mass and the corresponding target atom mass.  Candidate priorities are
carried by the registered common-atom marginal, while current and private
blockers enter through fresh-coordinate product laws, not through an old-cut
Boolean wrapper.

Then the atoms \(C^0_{\omega_0,\tau}\) and
\(C^{\prime0}_{\omega_0,\tau}\) are measurable; as \(\tau\) ranges over the
finite truth vectors, the source atoms form a finite pairwise-disjoint cover
of \(C^0_{\omega_0}\) and the target atoms form a finite pairwise-disjoint
cover of \(C^{\prime0}_{\omega_0}\); transported paired states lie in
corresponding truth-vector atoms; finite additivity gives
\[
\nu(C^0_{\omega_0})=\sum_{\tau\subseteq\Lambda}
  \nu(C^0_{\omega_0,\tau}),\qquad
\nu'(C^{\prime0}_{\omega_0})=\sum_{\tau\subseteq\Lambda}
  \nu'(C^{\prime0}_{\omega_0,\tau});
\]
and for each \((\omega_0,\tau)\) there is \(w_{\omega_0,\tau}\ge0\) with
\[
\operatorname{ofReal}(w_{\omega_0,\tau})
=\nu(C^0_{\omega_0,\tau})
=\nu'(C^{\prime0}_{\omega_0,\tau}).
\]
These are precisely the atom-mass and finite-additivity inputs needed by
\(\noderef{SamplingSplitOutsideBlockerStageCellCommonRefinement}\).
\end{helper}
\begin{proof}
Fix \(\omega_0\).  Measurability of a truth-vector atom follows from the
measurability of the base cell and the finitely many label sets: the atom is
the base cell intersected with, for each \(\lambda\), either \(Q_\lambda\) or
its complement.  The same finite-intersection argument applies on the target
side.

Because \(\Lambda\) is finite, every source state in \(C^0_{\omega_0}\) has
a unique truth vector
\[
\tau(\eta)=\{\lambda\in\Lambda:\eta\in Q_\lambda\}.
\]
This proves that the source truth-vector atoms cover the base cell.  If two
truth vectors differ, some label \(\lambda\) has different membership status
in them, so no state can lie in both corresponding atoms.  Hence the cover is
pairwise disjoint.  The target cover and target disjointness are identical
with \(Q'_\lambda\) in place of \(Q_\lambda\).

Now let \(T_{\omega_0}(\eta,\eta')\) hold.  The paired-cell transport says
\(\eta\in C^0_{\omega_0}\) iff
\(\eta'\in C^{\prime0}_{\omega_0}\), and label transport says
\(\eta\in Q_\lambda\) iff \(\eta'\in Q'_\lambda\) for every
\(\lambda\).  Combining these equivalences with the definition of the
truth-vector atoms gives
\[
\eta\in C^0_{\omega_0,\tau}
\Longleftrightarrow
\eta'\in C^{\prime0}_{\omega_0,\tau}.
\]

The mass comparison is the only point at which the split-blocker product
analysis enters this bookkeeping lemma.  It enters as a hypothesis of the
present statement: for each \((\omega_0,\tau)\) there is a nonnegative real
\(w_{\omega_0,\tau}\) whose extended nonnegative value is both the source
atom mass and the transported target atom mass.  This is exactly the
supported-common-atom output provided upstream by the actual positive-mass
product-subcell producer together with the zero-mass branch; the present node
does not derive that analytic product fact and does not repackage the fixed
stage calculation.

Finally apply finite additivity to the finite measurable disjoint covers just
constructed.  This gives the two displayed base-to-atom sum identities.  The
transported atom-mass identity and the nonnegative representatives supplied
by the atom-mass hypothesis are exactly the remaining inputs required by
\(\noderef{SamplingSplitOutsideBlockerStageCellCommonRefinement}\).
\end{proof}
